% !TEX root = ../../main.tex
% C6.2 — Hiện thực tiến trình xử lý nền AI (RÚT GỌN 2026-09-29; bản cũ ở report/archive/)
% Khớp backend: ai/registry.py (allowlist, SHA-256, giấy phép, PRODUCTION_PREFLIGHT_IDS, resolve() kiểm lại mã băm),
% config/models.example.json (4 detector, 2 tracker, 1 encoder), ai/configuration.py (ConfigApplyCoordinator: nạp thử trước khi áp dụng),
% ai/detectors/ultralytics.py + ai/encoders/image.py (mô hình chạy trong tiến trình con, quá hạn 120 s thì dừng), ai/encoders/rasa.py,
% config/ai_resources.json (worker chính thức chỉ áp dụng số luồng = 2; các ngưỡng RAM/đĩa/hàng đợi/điểm ảnh chỉ dùng trong preflight của worker demo nên KHÔNG đưa vào báo cáo),
% workers/production.py (xử lý từng khung, đóng ảnh ngay, kiểm tra hạn công việc mỗi khung), workers/production_main.py (mỗi công việc một tiến trình con, hạn 3600 s),
% workers/durable.py (khóa toàn cục, mã lỗi theo bước, thử lại 5 s nhân đôi tới 300 s tối đa 3 lần, RTSP không thử lại, dừng êm để lại RUNNING),
% services/diagnostics.py (kiểm tra AI: 6 khung hình, cách 5, tối đa 90 khung hình nguồn, hạn 180 s).
\section{Hiện thực tiến trình xử lý nền AI}
\label{sec:6.2}

Mục~\ref{sec:5.2.1} đã trình bày các bước xử lý và vòng đời của một công việc. Mục này trình bày cách tiến trình xử lý nền được hiện thực để chạy ổn định trên một máy chỉ có CPU.

\subsection{Quản lý và đóng gói mô hình}
\label{sec:6.2.1}

Hệ thống chỉ nạp những mô hình được khai báo trong một \emph{danh sách đăng ký mô hình}, mỗi mục ghi mã, phiên bản, tệp trọng số kèm mã băm SHA-256, phiên bản tiền xử lý và thông tin nguồn gốc, gồm cả giấy phép; với Tracker còn có danh sách Detector tương thích. Một mô hình chỉ dùng được khi tệp trọng số tồn tại, mã băm khớp, giấy phép đã được phê duyệt và mô hình đã được nạp thử thành công; mô hình không thỏa điều kiện vẫn được liệt kê trên màn hình \emph{Mô hình AI} với nhãn \emph{Không khả dụng}. Hiện YOLO11n, ByteTrack, BoT-SORT và RaSa dùng được, còn các phiên bản YOLO lớn hơn và YOLOX chưa có tệp trọng số (Mục~\ref{sec:6.1.4}). Mã băm được tính lại mỗi lần một cặp mô hình được chọn để chạy, nên một tệp trọng số bị thay trên đĩa không thể âm thầm được dùng. Khi Quản trị viên áp dụng cấu hình Detector-Tracker mới, máy chủ nạp thử cả ba thành phần trước khi lưu, và mỗi công việc dùng đúng phiên bản cấu hình tại thời điểm được tạo.

Mỗi loại mô hình được bọc trong một \emph{bộ chuyển đổi} (adapter) có cùng giao diện: mở để nạp mô hình, xử lý một đầu vào, và đóng để giải phóng tài nguyên. Bộ chuyển đổi kiểm tra đầu ra trước khi chuyển cho bước sau: Detector chỉ giữ lớp người và khung bao hợp lệ; vector của RaSa phải đủ 256 chiều, không chứa giá trị không hợp lệ và đã được chuẩn hóa độ dài. Nhờ vậy, thêm một mô hình mới, chẳng hạn kích hoạt YOLOX, chỉ cần thêm tệp trọng số và mục đăng ký, cùng bộ chuyển đổi nếu mô hình thuộc loại chưa có.

\subsection{Giới hạn tài nguyên và xử lý lỗi}
\label{sec:6.2.2}

Vì máy triển khai có bộ nhớ hạn chế, tiến trình xử lý nền được giới hạn ở nhiều mức. Số luồng tính toán của các thư viện số học được đặt bằng 2, để mô hình không chiếm hết các nhân CPU mà máy chủ ứng dụng và các kho dữ liệu cũng cần. Detector và bộ mã hóa ảnh chạy trong một tiến trình con riêng; nếu một lần suy luận không trả kết quả trong 120 giây, tiến trình con bị dừng hẳn, vì một lời gọi bị treo trong thư viện gốc không thể được ngắt từ trong cùng tiến trình. Mỗi công việc được chạy trong một tiến trình riêng với thời hạn tổng mặc định 3.600 giây; khi công việc kết thúc, tiến trình kết thúc theo, nên bộ nhớ của mô hình và khung hình được hệ điều hành thu hồi hoàn toàn. Khung hình được xử lý lần lượt từng khung và giải phóng ngay, trừ các ứng viên khung hình đại diện vốn đã bị giới hạn.

Mỗi lỗi được gắn với bước gây ra lỗi và một mã lỗi cụ thể, lưu cùng công việc để hiển thị trên màn hình quản trị, và được xếp trước vào nhóm \emph{tạm thời} hoặc \emph{không phục hồi được}. Với tệp video, lỗi tạm thời được thử lại sau 5 giây, nhân đôi sau mỗi lần và không quá 300 giây, tối đa ba lần. Phiên RTSP không được thử lại, vì đoạn video đã trôi qua; bộ lập lịch sẽ tạo phiên mới sau 5 phút (Mục~\ref{sec:5.2.1}). Khi tiến trình nhận tín hiệu dừng, công việc đang chạy được giữ ở trạng thái \emph{Đang xử lý} để được nhận lại khi thời hạn giữ chỗ hết, thay vì bị đánh dấu thất bại ngay. Một khóa toàn cục trong PostgreSQL bảo đảm chỉ một tiến trình xử lý công việc tại một thời điểm, kể cả khi tiến trình xử lý nền vô tình được khởi động hai lần.

\paragraph{Kiểm tra AI.}\label{sec:6.2.5} Màn hình \emph{Kiểm tra AI} chạy thử các mô hình trên dữ liệu thật mà không tạo công việc: với một camera, hệ thống đọc tối đa 90 khung hình từ luồng RTSP, lấy 6 khung hình rồi lần lượt chạy Detector, Tracker và bộ mã hóa ảnh; với nhóm tìm kiếm, hệ thống mã hóa một ảnh và một câu mô tả mẫu. Mỗi lần kiểm tra giới hạn trong 180 giây, và chỉ một lần kiểm tra được chạy tại một thời điểm.
